Silicon photonics has emerged as a leading platform for large-scale
photonic integration owing to its compatibility with complementary
metal-oxide-semiconductor (CMOS) fabrication infrastructure, high
index contrast, and the availability of multi-project wafer (MPW)
services~\cite{reed2010silicon, bogaerts2012silicon}.  Among the
building blocks of silicon photonic integrated circuits (PICs), microring
resonators (MRRs) and their elongated variants---racetrack
resonators---play a pivotal role in wavelength-selective routing,
filtering, modulation, and sensing~\cite{little1997microring,
yariv2000critical}.

Dense wavelength division multiplexing (DWDM) leverages the high
spectral efficiency of closely spaced optical channels to maximise
the information capacity of a single fibre or waveguide
bus~\cite{winzer2012optical}. In the C-band,
the International Telecommunication Union (ITU) G.694.1 standard
defines a \SI{100}{\giga\hertz} (\SI{\sim 0.8}{\nano\meter}) channel
grid, placing stringent requirements on the resonance wavelength
accuracy and thermal stability of ring-based demultiplexers.

Thermal tuning via integrated microheaters is the most widely adopted
post-fabrication technique for compensating fabrication-induced
resonance shifts and for dynamic channel
assignment~\cite{dong2010thermally, watts2013adiabatic}. Titanium
nitride (TiN) is preferred as the heater material in CMOS-compatible
flows because of its low electrical resistivity, mechanical stability,
and optical absorption properties that minimise parasitic waveguide
loss~\cite{sacher2014wide}.

Despite the broad literature on individual ring modulators and switches,
a complete open-source workflow---from parametric GDSII layout generation
through electromagnetic simulation to manuscript-ready figures---for a
multi-channel DWDM demultiplexer based on racetrack MRRs is still
lacking.  This paper addresses that gap with the following contributions:

\begin{enumerate}
  \item A parametric gdsfactory-based design script that generates
        foundry-ready GDSII layouts for a configurable number of ITU
        C-band channels (Section~\ref{sec:layout}).
  \item Full-wave FDTD simulation using MEEP of the through- and
        drop-port spectral responses, including temperature-dependent
        refractive index shifts (Section~\ref{sec:simulation}).
  \item Systematic extraction of key performance metrics: extinction
        ratio (ER), insertion loss (IL), quality factor $Q$, free
        spectral range (FSR), and inter-channel crosstalk
        (Section~\ref{sec:results}).
  \item Characterisation of the thermal tuning efficiency and power
        budget for a \SI{\pm 2}{\nano\meter} tuning range
        (Section~\ref{sec:thermal}).
\end{enumerate}

The remainder of the paper is organised as follows.
Section~\ref{sec:design} introduces the device architecture and
physical parameters.  Section~\ref{sec:layout} describes the automated
layout generation flow.  Sections~\ref{sec:simulation}--\ref{sec:thermal}
present the simulation methodology and results.
Section~\ref{sec:conclusion} concludes the paper.
